% !TEX root = ../main.tex
\chapter{Thiết lập thực nghiệm và đánh giá hệ thống}
\label{Chapter4}

\par\noindent\rule{\textwidth}{0.4pt}\par

\section{Mục tiêu và phạm vi đánh giá}

Chương 3 đã phân chia ConferenceSpace thành ba lớp trách nhiệm: nghiệp vụ cốt lõi, thuật toán xác định có thể kiểm chứng và AI hỗ trợ. Chương này đánh giá mức độ hoàn thành trách nhiệm của từng lớp trong các điều kiện thử nghiệm đã xác định, sau đó đối chiếu kết quả kỹ thuật với phản hồi của người dùng. Do mỗi lớp tạo ra một loại đầu ra khác nhau, quy trình kiểm thử sử dụng các nguồn bằng chứng riêng thay vì quy toàn bộ hệ thống về một thước đo chung.

Bảng~\ref{tab:ch4-demo-questions} trình bày bốn câu hỏi đánh giá, nguồn bằng chứng và phạm vi kết luận tương ứng.

\begingroup
\setlength{\tabcolsep}{4pt}
\renewcommand{\arraystretch}{1.3}
\begin{longtable}{|>{\raggedright\arraybackslash}p{0.26\textwidth}|>{\raggedright\arraybackslash}p{0.26\textwidth}|>{\raggedright\arraybackslash}p{0.22\textwidth}|>{\raggedright\arraybackslash}p{0.18\textwidth}|}
\caption{Câu hỏi đánh giá và phạm vi bằng chứng}
\label{tab:ch4-demo-questions} \\
\hline
\TableHeader{Câu hỏi đánh giá} & \TableHeader{Nguồn bằng chứng} & \TableHeader{Chỉ số chính} & \TableHeader{Phạm vi kết luận} \\ \hline
\endfirsthead
\hline
\TableHeader{Câu hỏi đánh giá} & \TableHeader{Nguồn bằng chứng} & \TableHeader{Chỉ số chính} & \TableHeader{Phạm vi kết luận} \\ \hline
\endhead
\hline
\multicolumn{4}{|r|}{\textit{(Tiếp tục ở trang sau...)}} \\ \hline
\endfoot
\hline
\endlastfoot
Hệ thống nghiệp vụ có đủ nhanh và ổn định trong điều kiện tải thử nghiệm hay không? & Kiểm thử tải bằng k6, thống kê yêu cầu thất bại và giám sát tài nguyên & Thông lượng, độ trễ p95, tỷ lệ yêu cầu thất bại, mức sử dụng CPU và RAM & Các đường xử lý và cấu hình đã chọn \\ \hline
Thuật toán đối sánh và phát hiện xung đột lợi ích có chạy nhanh, cho kết quả có thể tái lập và có thể giải thích hay không? & Go microbenchmark và thử nghiệm trên dữ liệu tổng hợp mô phỏng lược đồ Semantic Scholar & Thời gian xử lý, bộ nhớ, Hit@k, MRR, nDCG, độ phủ và độ đồng đều của phân công & Chi phí mã thuật toán và hành vi trên dữ liệu tổng hợp \\ \hline
Các luồng AI có tạo đầu ra đủ bám bằng chứng và giữ đúng ranh giới hỗ trợ hay không? & Đối chiếu trực tiếp, Textual Claim-based Assessment (TCA) và thử nghiệm hội thoại & Exact Match, ROUGE, F1, Truthfulness, Coverage, Additionality và kết quả hội thoại & Từng tác vụ và từng bộ kiểm thử tương ứng \\ \hline
Kết quả kỹ thuật có tương ứng với trải nghiệm của người dùng hay không? & Khảo sát sau sử dụng theo ba vai trò & Mức hài lòng, cảm nhận giảm thao tác và phản hồi định tính & Cảm nhận của mẫu người dùng tham gia \\
\end{longtable}
\endgroup

Hiệu năng của backend không chứng minh đầu ra AI đáng tin cậy; mức bám bằng chứng cao không chứng minh quyết định học thuật đúng; phản hồi tích cực của người dùng cũng không thay thế kiểm thử định lượng. Chương này vì vậy diễn giải từng kết quả trong phạm vi của phép đo đã tạo ra kết quả đó.

\subsection{Chuỗi thực nghiệm và ranh giới với Phụ lục C}

Chuỗi thực nghiệm được tổ chức theo quan hệ từ đầu vào đến kết luận. Nhóm khởi tạo dữ liệu và điều kiện chạy, thực hiện thao tác tương ứng với từng kịch bản, lưu đầu ra cùng thông tin truy vết, tính các chỉ số đã xác định, rồi phân tích kết quả trong phạm vi của phép đo. Hình~\ref{fig:ch4-evaluation-chain} minh họa chuỗi bằng chứng ở mức khái quát.

\begin{figure}[H]
  \centering
  \begin{adjustbox}{max width=\textwidth}
  \begin{tikzpicture}[
    node distance=0.55cm and 0.35cm,
    evalstep/.style={draw, rounded corners, align=center, minimum height=0.85cm, text width=2.25cm, fill=gray!8},
    evalarrow/.style={-{Latex[length=2mm]}, thick}
  ]
    \node[evalstep] (input) {Đầu vào và\\ điều kiện};
    \node[evalstep, right=of input] (run) {Thao tác theo\\ kịch bản};
    \node[evalstep, right=of run] (record) {Đầu ra và\\ tư liệu truy vết};
    \node[evalstep, right=of record] (metric) {Chỉ số và\\ đối chiếu};
    \node[evalstep, right=of metric] (claim) {Phân tích và\\ giới hạn kết luận};
    \draw[evalarrow] (input) -- (run);
    \draw[evalarrow] (run) -- (record);
    \draw[evalarrow] (record) -- (metric);
    \draw[evalarrow] (metric) -- (claim);
  \end{tikzpicture}
  \end{adjustbox}
  \caption{Chuỗi thực nghiệm từ đầu vào đến kết luận}
  \label{fig:ch4-evaluation-chain}
\end{figure}

Chương 4 giữ phần thông tin cần thiết để người đọc hiểu mục tiêu, đối tượng, đầu vào, kịch bản, thao tác, giá trị của phép đo và giới hạn diễn giải. Phụ lục C không thay thế phần trình bày này; phụ lục cung cấp mức chi tiết sâu hơn về điều kiện khởi tạo, bộ dữ liệu kiểm thử, bản dữ liệu cố định, đơn vị đo, định dạng tư liệu và các khoảng trống truy vết. Vì vậy, người đọc có thể đọc Chương 4 để nắm mạch đánh giá, sau đó sử dụng các mục~\ref{appc:index}--\ref{appc:operations} để kiểm tra hoặc tái lập từng nhóm phép đo.

\section{Thiết lập thực nghiệm}

\subsection{Dữ liệu và nhóm thực nghiệm}

Đối với lớp AI, nhóm chọn một tập con của ReviewRebuttal (Re2), bộ dữ liệu chứa hồ sơ phản biện bài báo qua nhiều giai đoạn từ OpenReview~\cite{ch4ref1}. Tập dữ liệu gồm các bản ghi có đủ thành phần cần thiết để thực hiện lại ít nhất một luồng được đánh giá, chẳng hạn bài nộp, siêu dữ liệu, nhận xét phản biện, phản hồi của tác giả hoặc bản tổng hợp nhận xét. Bộ thực thi đã xử lý thành công 1.097 trên 1.127 bài đầu vào. Ba mươi trường hợp không hoàn tất do lỗi dịch vụ AI hoặc lỗi trong quá trình chạy benchmark, chiếm 2,66\% tổng số bài. Tỷ lệ hoàn tất 97,34\% cho thấy quy trình chạy ổn định trên phần lớn tập dữ liệu đã chọn. Các bài không hoàn tất không được đưa vào phép chấm TCA vì không tạo ra kết quả đủ điều kiện, nhưng vẫn được giữ trong thống kê vận hành để phản ánh đầy đủ tỷ lệ thành công của bộ thực thi.

Hình~\ref{fig:ch4-demo-dataset-distribution} trình bày phân bố 1.127 bài theo hội nghị hoặc chuyên đề. Bốn nguồn lớn nhất chiếm 70,72\% tập dữ liệu; do đó, kết quả cần được đọc như bằng chứng trên tập đã chọn, không phải phân bố đại diện cho mọi hội nghị.

\begin{figure}[!htbp]
  \centering
  \includegraphics[width=0.90\textwidth]{fig01b_runner_by_conference.png}
  \caption{Phân bố bài đầu vào theo hội nghị và chuyên đề}
  \label{fig:ch4-demo-dataset-distribution}
\end{figure}

Các nhóm thực nghiệm còn lại sử dụng tập kiểm thử phù hợp với từng trách nhiệm. Kiểm thử backend dùng dữ liệu tổng hợp gồm 300 hội nghị, 15.000 bài nộp và 9.000 quan hệ phản biện viên--hội nghị. Đánh giá đối sánh dùng 60 hồ sơ tác giả và 2.565 bài báo tổng hợp. Tự động điền bài nộp (Submission Autofill) có thêm 48 trường hợp kiểm thử gợi ý chuyên đề; Kiểm tra trước khi nộp bài (Submission Gating) có 8 trường hợp kiểm thử luật và 24 trường hợp kiểm thử cảnh báo nội dung; Chatbot Agent có 40 hội thoại thuộc tám nhóm kịch bản. Phụ lục~\ref{AppendixC} trình bày phân bố dữ liệu, điều kiện chạy và quy trình chấm chi tiết.

\subsection{Kịch bản và quy tắc đánh giá}

Mỗi nhóm đánh giá trả lời một câu hỏi riêng. Nhóm hiệu năng backend kiểm tra ba đường xử lý HTTP trên dữ liệu tổng hợp; nhóm thuật toán sử dụng bộ kiểm thử Go ở nhiều quy mô và một bản dữ liệu cố định tổng hợp để đánh giá chi phí xử lý, hành vi xếp hạng và kết quả phân công. Cấu hình khởi tạo và quy tắc tạo dữ liệu được trình bày tại các Mục~\ref{appc:backend-contract} và~\ref{appc:algorithm-contract}.

Ba kịch bản backend gồm truy vấn đọc, gợi ý phản biện viên và kiểm tra xung đột lợi ích. Kịch bản truy vấn đọc lần lượt liệt kê hội nghị, liệt kê bài nộp và tìm kiếm người dùng; mặc dù được gọi là CRUD trong tư liệu tải, kịch bản chỉ thực hiện yêu cầu đọc. Kịch bản gợi ý phản biện viên gọi đường xử lý \CodeBreak{reviewer-suggestions} nhưng không tự động phân công. Kịch bản xung đột lợi ích gọi đường xử lý kiểm tra xung đột, trong đó có truy vấn mẫu quan hệ đồng tác giả trên Neo4j bằng Cypher~\cite{ch4ref6}.

Đối với xếp hạng, nhóm giữ lại một bài từ mỗi hồ sơ tác giả làm truy vấn, loại bài đó khỏi hồ sơ tương ứng rồi xếp hạng các hồ sơ còn lại bằng Jaccard, \CodeBreak{overlap\_count} và phương pháp ngẫu nhiên. Kết quả được mô tả bằng Hit@k, MRR và nDCG@10, là các chỉ số thường dùng để đánh giá danh sách xếp hạng~\cite{ch4ref8}. Đối với phân công, nhóm so sánh Greedy với gán tuần tự và ngẫu nhiên bằng độ phủ, độ đồng đều, điểm Jaccard, lượt dự phòng và số ca tự gán. Vì chưa có phương án phân công tối ưu do Chủ tọa xác nhận, nhãn tác giả gốc và các chỉ số nội tại chỉ là bằng chứng gián tiếp.

Các luồng AI thực hiện những trách nhiệm khác nhau nên không được đánh giá bằng một chỉ số chung. Đối với đầu ra có dữ liệu tham chiếu xác định, nhóm đối chiếu trực tiếp kết quả của hệ thống với giá trị kỳ vọng. Đối với đầu ra văn bản mở, nhóm sử dụng TCA để kiểm tra quan hệ giữa từng mệnh đề và nguồn bằng chứng. Chatbot Agent được đánh giá riêng theo kịch bản hội thoại, quyền truy cập và hành vi gọi công cụ. Bảng~\ref{tab:ch4-ai-evaluation-contracts} tóm tắt đầu vào, đầu ra và cách kiểm tra của từng luồng; Phụ lục~\ref{AppendixC} trình bày cấu trúc dữ liệu, trường hợp kiểm thử và quy trình tái lập.

\begingroup
\small
\setlength{\tabcolsep}{2pt}
\renewcommand{\arraystretch}{1.2}
\begin{longtable}{|>{\RaggedRight\arraybackslash\hyphenpenalty=10000\exhyphenpenalty=10000}p{0.16\textwidth}|>{\RaggedRight\arraybackslash}p{0.16\textwidth}|>{\RaggedRight\arraybackslash}p{0.17\textwidth}|>{\RaggedRight\arraybackslash}p{0.24\textwidth}|>{\RaggedRight\arraybackslash}p{0.18\textwidth}|}
\caption{Đầu vào, đầu ra và cách đánh giá các luồng AI hỗ trợ}
\label{tab:ch4-ai-evaluation-contracts} \\
\hline
\TableHeader{Luồng} & \TableHeader{Đầu vào} & \TableHeader{Đầu ra quan sát} & \TableHeader{Cách kiểm tra} & \TableHeader{Phạm vi kết luận} \\ \hline
\endfirsthead
\hline
\TableHeader{Luồng} & \TableHeader{Đầu vào} & \TableHeader{Đầu ra quan sát} & \TableHeader{Cách kiểm tra} & \TableHeader{Phạm vi kết luận} \\ \hline
\endhead
\hline
\multicolumn{5}{|r|}{\textit{(Tiếp tục ở trang sau...)}} \\ \hline
\endfoot
\hline
\endlastfoot
Submission Autofill & Bản thảo, siêu dữ liệu tham chiếu và danh sách chuyên đề hợp lệ & Tiêu đề, tác giả, từ khóa, tóm tắt, trường bắt buộc và chuyên đề được gợi ý & Đối chiếu siêu dữ liệu bằng \mbox{Exact Match}, F1 và ROUGE; kiểm tra tính hợp lệ của chuyên đề được gợi ý & Độ khớp trường và khả năng hoàn tất; chưa xác nhận chuyên đề phù hợp nhất \\ \hline
Submission Gating: kiểm tra luật & Tệp bài nộp, siêu dữ liệu và quy tắc hội nghị & Phán quyết và mã quy tắc & Đối chiếu với kết quả kỳ vọng của từng trường hợp kiểm thử & Việc áp dụng các quy tắc đã định nghĩa \\ \hline
Submission Gating: cảnh báo nội dung & Nội dung bản thảo và phạm vi cảnh báo & Danh sách cảnh báo không chặn & Kiểm tra khả năng hoàn tất, số cảnh báo và việc không tự động chặn & Hành vi hỗ trợ; chưa xác nhận độ chính xác của từng cảnh báo \\ \hline
Reviewer Initial Analysis & Bản thảo và siêu dữ liệu bài nộp & Trích dẫn, chú giải và điểm cần lưu ý & Đối chiếu trích dẫn với bản thảo; dùng TCA cho các nhận định mở và nội dung tham chiếu & Mức bám nguồn và quan hệ với tham chiếu; không thay thế phán đoán của Phản biện viên \\ \hline
Review Quality Auditor & Bản phản biện và tiêu chí kiểm tra chất lượng & Phát hiện, trạng thái và mức độ nghiêm trọng & Dùng TCA để kiểm tra mức bám nguồn, tính hợp lệ và phần đồng thời đạt hai điều kiện & Chất lượng phát hiện; trạng thái \textit{block} không phải quyết định từ chối thao tác \\ \hline
Chair Decision Copilot & Bài nộp, bản phản biện, phản hồi của tác giả và nội dung thảo luận khả dụng & Cơ sở bằng chứng, điểm đồng thuận, điểm bất đồng và bản tổng hợp & Dùng TCA để kiểm tra mức bám nguồn và quan hệ với nội dung tham chiếu & Chất lượng bản tổng hợp; không đo độ chính xác của quyết định học thuật \\ \hline
Chatbot Agent & Vai trò, yêu cầu của người dùng và dữ liệu được phép truy cập & Câu trả lời, chuỗi gọi công cụ và trạng thái hội thoại & Rà soát thủ công theo kịch bản, công cụ dự kiến và tiêu chí hành vi & Hành vi trong 40 hội thoại đã thiết kế \\
\end{longtable}
\endgroup

Đối với các luồng AI, bộ thực thi tạo và lưu đầu ra trước khi hoạt động chấm bắt đầu. Submission Autofill và nhánh kiểm tra luật của Submission Gating được đối chiếu trực tiếp; Reviewer Initial Analysis, Review Quality Auditor và Chair Decision Copilot được đánh giá theo quan hệ mệnh đề--bằng chứng; Chatbot Agent được chấm theo hành vi hội thoại. Hình~\ref{fig:ch4-ai-evaluation-flow} minh họa ba nhánh đánh giá này.

\begin{figure}[htbp]
  \centering
  \begin{adjustbox}{max width=0.88\textwidth}
  \begin{tikzpicture}[>=Latex,
    block/.style={draw=black!65, rounded corners, align=center, text width=2.9cm, minimum height=0.85cm, fill=gray!8},
    arrow/.style={-{Latex[length=2mm]}, thick}
  ]
    \node[block] (data) at (0,3.6) {Bài nộp và\\ ngữ cảnh nguồn};
    \node[block] (worker) at (0,1.4) {Bộ thực thi\\ luồng xử lý};
    \node[block] (outputs) at (0,-0.8) {Gói đầu ra\\ và điểm kiểm tra};
    \node[block] (direct) at (4.2,-3.3) {Đối chiếu trực tiếp\\ Autofill / Gating};
    \node[block] (tca) at (-4.2,-3.3) {TCA cho\\ đầu ra văn bản};
    \node[block] (chat) at (0,-3.3) {Chấm theo kịch bản\\ Chatbot Agent};
    \draw[arrow] (data) -- (worker);
    \draw[arrow] (worker) -- (outputs);
    \draw[arrow] (outputs) -- (direct);
    \draw[arrow] (outputs) -- (tca);
    \draw[arrow] (outputs) -- (chat);
  \end{tikzpicture}
  \end{adjustbox}
\caption{Các nhánh đánh giá đầu ra của luồng xử lý AI}
\label{fig:ch4-ai-evaluation-flow}
\end{figure}

\subsection{Môi trường và tư liệu thực nghiệm}

Bảng~\ref{tab:ch4-demo-environments} liên kết từng nhóm đánh giá với môi trường thực hiện, quy mô chính và tư liệu đầu ra. Bảng chỉ cung cấp thông tin cần thiết để xác định nguồn của phép đo; cấu hình chi tiết, đơn vị đo và yêu cầu truy vết được trình bày trong Phụ lục~\ref{AppendixC}.

\begingroup
\setlength{\tabcolsep}{2.5pt}
\renewcommand{\arraystretch}{1.3}

\begin{longtable}{|>{\raggedright\arraybackslash}p{0.17\textwidth}|>{\raggedright\arraybackslash}p{0.40\textwidth}|>{\raggedright\arraybackslash}p{0.18\textwidth}|>{\raggedright\arraybackslash}p{0.20\textwidth}|}
\caption{Môi trường và tư liệu thực nghiệm}
\label{tab:ch4-demo-environments} \\
\hline
\TableHeader{Nhóm thực nghiệm} & \TableHeader{Môi trường hoặc phương pháp} & \TableHeader{Quy mô chính} & \TableHeader{Tư liệu đầu ra} \\ \hline
\endfirsthead
\hline
\TableHeader{Nhóm thực nghiệm} & \TableHeader{Môi trường hoặc phương pháp} & \TableHeader{Quy mô chính} & \TableHeader{Tư liệu đầu ra} \\ \hline
\endhead
\hline
\multicolumn{4}{|r|}{\textit{(Tiếp tục ở trang sau...)}} \\ \hline
\endfoot
\hline
\endlastfoot
Hiệu năng backend & k6 tạo tải HTTP lên API; yêu cầu được xử lý qua PostgreSQL, Neo4j và Redis; 20 người dùng ảo trong 30 giây cho mỗi kịch bản & 300 hội nghị; 15.000 bài; 9.000 quan hệ phản biện viên--hội nghị & Tóm tắt k6 và nhật ký tài nguyên \\ \hline
Chi phí thuật toán & Go microbenchmark chạy trực tiếp trên mã và lặp 5 lần; không gồm HTTP, cơ sở dữ liệu và mạng & Nhiều quy mô đầu vào & Thời gian, bộ nhớ và số lần cấp phát trên mỗi thao tác \\ \hline
Xếp hạng và phân công & Leave-one-out theo tác giả; so sánh Jaccard, đếm số chủ đề chung, gán tuần tự và ngẫu nhiên & 60 hồ sơ; 2.565 bài tổng hợp & Báo cáo Markdown/CSV và bảng chỉ số \\ \hline
Luồng AI & Bộ điều phối và các tiến trình xử lý tạo đầu ra cho từng bài; Gemini 3.1 Flash Lite được gọi qua bộ định tuyến mô hình & 1.127 bài & Ngữ cảnh nguồn, đầu ra, trạng thái, thời gian và token \\ \hline
TCA & \CodeBreak{dleemiller/ModernCE-large-nli} phân loại quan hệ bằng chứng--mệnh đề; Qwen3.5-2B và Qwen3.5-0.8B thực hiện chuẩn hóa và phân loại mức độ hỗ trợ~\cite{ch4ref2,ch4ref3} & 1.097 bài có kết quả đủ điều kiện & JSONL theo bài và nhóm chỉ số \\ \hline
Chatbot Agent & Thử nghiệm thủ công hồi cứu theo vai trò, công cụ dự kiến được gọi và tiêu chí định trước & 40 hội thoại & Bản ghi hội thoại và nhật ký gọi công cụ \\ \hline
Khảo sát sau sử dụng & Bảng hỏi theo vai trò sau khi người tham gia trải nghiệm hệ thống & 91 phản hồi & Điểm số và phản hồi định tính \\
\end{longtable}
\endgroup

\subsection{Phương pháp đánh giá theo nhóm}

Kiểm thử tải dùng k6 để tạo tải HTTP lên API và thu các chỉ số tích hợp về số yêu cầu, thời gian xử lý và tỷ lệ thất bại~\cite{ch4ref4}. Go microbenchmark chạy trực tiếp trên mã thuật toán theo cơ chế benchmark của gói \CodeBreak{testing}~\cite{ch4ref5}, nhằm tách thời gian xử lý và mức sử dụng bộ nhớ khỏi chi phí HTTP, cơ sở dữ liệu và mạng. Phép thử xếp hạng dùng phương pháp leave-one-out theo tác giả; phép thử phân công so sánh Greedy với gán tuần tự và ngẫu nhiên bằng các chỉ số nội tại vì chưa có phương án phân công tối ưu do Chủ tọa xác nhận.

Các đầu ra AI được đánh giá theo ba cách. Tự động điền bài nộp (Submission Autofill) được đối chiếu với siêu dữ liệu tham chiếu bằng Exact Match, F1 và ROUGE, trong đó ROUGE đo mức trùng lặp giữa văn bản tạo ra và văn bản tham chiếu~\cite{ch4ref9}; phần kiểm tra luật của Submission Gating được đối chiếu với phán quyết và mã luật dự kiến. Phân tích và tóm tắt bài báo (Reviewer Initial Analysis), Kiểm tra chất lượng nhận xét (Review Quality Auditor) và Báo cáo tổng hợp hỗ trợ Chủ tọa (Chair Decision Copilot) được đánh giá bằng TCA. Chatbot Agent được kiểm tra qua các kịch bản hội thoại theo vai trò, quyền truy cập và phạm vi hỗ trợ.

\subsection{Textual Claim-based Assessment}
\label{subsec:ch4-tca}

Các đầu ra như bản phân tích dành cho Phản biện viên, phát hiện về chất lượng nhận xét và bản tổng hợp dành cho Chủ tọa không có một đáp án văn bản duy nhất. Các phép đo dựa trên mức trùng khớp từ vựng có thể đánh giá thấp một cách diễn đạt khác nhưng đúng, đồng thời không xác định được một nhận định có thực sự được nguồn dữ liệu hỗ trợ hay không. Mặt khác, nội dung do con người tạo không nhất thiết bao quát toàn bộ thông tin có căn cứ trong bài báo. Vì vậy, nhóm sử dụng Textual Claim-based Assessment (TCA), một khung hậu kiểm tự động đánh giá đầu ra văn bản theo đơn vị mệnh đề--chứng cứ.

TCA không tạo lại đầu ra của các luồng AI. Sau khi bộ thực thi hoàn tất một lượt chạy, hệ thống lưu đầu ra cùng ngữ cảnh nguồn, nội dung tham chiếu và thông tin truy vết thành một gói kết quả. TCA chỉ đọc các gói đã lưu và chấm chúng theo cùng một quy trình. Cách tổ chức này tách hoạt động tạo nội dung khỏi hoạt động đánh giá, đồng thời cho phép kiểm tra lại kết quả mà không làm thay đổi đầu ra ban đầu.

Bộ thực thi tiếp nhận 1.127 bài và tạo được 1.097 gói kết quả đủ điều kiện cho TCA. Ba mươi lượt không hoàn tất vẫn được tính khi báo cáo tỷ lệ thành công của bộ thực thi, nhưng không được đưa vào phép chấm chất lượng vì thiếu đầu ra hợp lệ. Do đó, các tỷ lệ hoàn tất vận hành sử dụng 1.127 bài làm mẫu số, còn các chỉ số TCA được tổng hợp trên 1.097 bài có gói kết quả đủ điều kiện.

\subsubsection{Đầu vào và đầu ra}

Mỗi gói đầu vào của TCA gồm ba nhóm thông tin. Nhóm thứ nhất là đầu ra cần đánh giá, chẳng hạn điểm cần lưu ý, phát hiện về bản phản biện hoặc cơ sở bằng chứng dành cho Chủ tọa. Nhóm thứ hai là nguồn dùng để kiểm tra các nhận định, gồm bản thảo, bản phản biện, phản hồi của tác giả hoặc nội dung thảo luận tương ứng với từng luồng. Nhóm thứ ba là nội dung do con người tạo dùng làm tham chiếu khi phép đo cần xác định mức độ bao phủ hoặc phần nội dung bổ sung.

Đầu ra của TCA gồm kết quả kiểm tra ở cấp mệnh đề, các chỉ số tổng hợp theo bài và thông tin truy vết về nhánh đánh giá đã sử dụng. Các gói kết quả này được lưu ở định dạng JSONL để nhóm có thể tổng hợp số liệu mà vẫn truy ngược được từ chỉ số về bài và mệnh đề tương ứng. Mục~\ref{appc:tca} trình bày cấu trúc dữ liệu, cấu hình và quy tắc tổng hợp chi tiết.

\subsubsection{Các giai đoạn xử lý}

Hình~\ref{fig:ch4-tca-flow} trình bày bốn giai đoạn xử lý chung của TCA. Mỗi luồng sử dụng nguồn bằng chứng và nhóm chỉ số riêng, nhưng đều tuân theo chuỗi từ kiểm tra gói kết quả đến lưu kết quả có thể truy vết.

\begin{figure}[H]
  \centering
  \begin{adjustbox}{max width=\textwidth}
  \begin{tikzpicture}[
    node distance=0.35cm,
    tcastep/.style={draw, rounded corners, align=center, minimum height=1.1cm, text width=2.45cm, fill=gray!8},
    tcaarrow/.style={-{Latex[length=2mm]}, thick}
  ]
    \node[tcastep] (package) {Gói kết quả\\ đã lưu};
    \node[tcastep, right=of package] (eligibility) {Kiểm tra điều kiện\\ và chọn nhánh};
    \node[tcastep, right=of eligibility] (claims) {Tách và chuẩn hóa\\ đơn vị đánh giá};
    \node[tcastep, right=of claims] (evidence) {Đối chiếu với\\ nguồn bằng chứng};
    \node[tcastep, right=of evidence] (aggregate) {Tổng hợp chỉ số\\ và lưu truy vết};
    \draw[tcaarrow] (package) -- (eligibility);
    \draw[tcaarrow] (eligibility) -- (claims);
    \draw[tcaarrow] (claims) -- (evidence);
    \draw[tcaarrow] (evidence) -- (aggregate);
  \end{tikzpicture}
  \end{adjustbox}
  \caption{Luồng xử lý tổng quát của TCA}
  \label{fig:ch4-tca-flow}
\end{figure}

\textbf{Giai đoạn 1 -- kiểm tra điều kiện và chọn nhánh đánh giá.} TCA kiểm tra gói kết quả có đủ đầu ra và nguồn đối chiếu cần thiết hay không, sau đó chọn nhánh xử lý phù hợp với từng luồng. Các gói thiếu dữ liệu bắt buộc được ghi nhận bằng cờ chất lượng và không được đưa vào mẫu số của phép chấm tương ứng.

\textbf{Giai đoạn 2 -- tách và chuẩn hóa đơn vị đánh giá.} TCA chuyển đầu ra thành các đơn vị có thể kiểm tra độc lập, chẳng hạn trích dẫn, điểm cần lưu ý, phát hiện về bản phản biện hoặc nhận định trong bản tổng hợp. Quy trình đồng thời loại các câu thuần túy hành chính nếu chúng không chứa nhận định cần đối chiếu với nội dung học thuật.

\textbf{Giai đoạn 3 -- đối chiếu với nguồn bằng chứng.} Đối với trích dẫn trực tiếp, TCA kiểm tra mức khớp với văn bản nguồn. Đối với nhận định mở, TCA xác định quan hệ giữa mệnh đề và phần bằng chứng liên quan. Khi có nội dung tham chiếu do con người tạo, TCA tiếp tục kiểm tra mệnh đề nào trùng với tham chiếu và mệnh đề nào có căn cứ nhưng nằm ngoài tham chiếu.

\textbf{Giai đoạn 4 -- tổng hợp và lưu kết quả.} TCA tổng hợp kết quả ở cấp mệnh đề thành các chỉ số theo bài và từng luồng, đồng thời giữ lại mã nhận diện, cờ chất lượng và thông tin cần thiết để truy vết. Công thức, quy tắc xử lý trường hợp đặc biệt và cấu hình kỹ thuật được trình bày tại Mục~\ref{appc:tca}.

\subsubsection{Ý nghĩa và giới hạn của các chỉ số}

\begingroup
\small
\setlength{\tabcolsep}{3pt}
\renewcommand{\arraystretch}{1.2}
\begin{longtable}{|>{\RaggedRight\arraybackslash}p{0.19\textwidth}|>{\RaggedRight\arraybackslash}p{0.39\textwidth}|>{\RaggedRight\arraybackslash}p{0.34\textwidth}|}
\caption{Ý nghĩa và giới hạn diễn giải của các chỉ số TCA}
\label{tab:ch4-tca-metric-meanings} \\
\hline
\TableHeader{Chỉ số} & \TableHeader{Ý nghĩa trong phép đánh giá} & \TableHeader{Không được diễn giải thành} \\ \hline
\endfirsthead
\hline
\TableHeader{Chỉ số} & \TableHeader{Ý nghĩa trong phép đánh giá} & \TableHeader{Không được diễn giải thành} \\ \hline
\endhead
\hline
\multicolumn{3}{|r|}{\textit{(Tiếp tục ở trang sau...)}} \\ \hline
\endfoot
\hline
\endlastfoot
Tỷ lệ trích dẫn bám nguồn & Tỷ lệ trích dẫn có thể đối chiếu trực tiếp hoặc một phần với văn bản nguồn & Độ đúng của toàn bộ bản phân tích \\ \hline
Truthfulness & Mức độ các mệnh đề được nguồn dữ liệu chỉ định hỗ trợ & Độ chính xác của quyết định học thuật \\ \hline
Coverage & Mức độ đầu ra bao phủ các điểm có trong nội dung tham chiếu do con người tạo & Mức độ bao phủ toàn bộ nội dung bài báo \\ \hline
Additionality & Tỷ lệ nội dung được đánh giá là có căn cứ nhưng không trùng trực tiếp với tham chiếu của con người & Tính mới, mức hữu ích hoặc giá trị chuyên môn \\ \hline
Validity & Tỷ lệ phát hiện phù hợp với tiêu chí kiểm tra tương ứng & Bằng chứng rằng phát hiện đã bám nguồn đầy đủ \\ \hline
Tỷ lệ vừa bám nguồn vừa hợp lệ & Tỷ lệ phát hiện đồng thời có căn cứ và phù hợp với tiêu chí kiểm tra & Tỷ lệ hệ thống nên tự động chặn người dùng \\ \hline
Tỷ lệ rủi ro cao & Tỷ lệ bài có dấu hiệu cần ưu tiên kiểm tra lại theo quy tắc TCA & Tỷ lệ bài nên bị từ chối \\
\end{longtable}
\endgroup

Kết quả TCA được dùng làm bằng chứng gián tiếp về mức bám nguồn và quan hệ giữa đầu ra AI với nội dung tham chiếu. Phép chấm chưa được hiệu chuẩn bằng một bộ nhãn chuyên gia độc lập; vì vậy, kết quả không thay thế đánh giá của Phản biện viên hoặc Chủ tọa. Báo cáo chỉ sử dụng kết quả TCA để mô tả mức độ đầu ra có thể được truy vết về nguồn, phần nội dung trùng với tham chiếu và phần nội dung có căn cứ nằm ngoài tham chiếu. Các kết quả này không trực tiếp đo mức hữu ích chuyên môn, tính đầy đủ của nhận xét hoặc độ chính xác của quyết định chấp nhận bài.

\section{Đánh giá lớp nghiệp vụ cốt lõi}
\label{sec:chap4-core-business}

\subsection{Hiệu năng backend}

Hình~\ref{fig:ch4-k6-throughput} và Hình~\ref{fig:ch4-k6-latency} minh họa hai khía cạnh khác nhau của cùng phép đo: thông lượng cho biết số yêu cầu được xử lý trong một đơn vị thời gian, còn các phân vị độ trễ cho biết mức phản hồi trong từng kịch bản. Các hình không mở rộng phạm vi đo vượt quá ba đường xử lý HTTP đã mô tả.

\begin{figure}[!htbp]
  \centering
  \includegraphics[width=0.95\textwidth]{fig02_k6_throughput.png}
  \caption{Thông lượng HTTP theo kịch bản k6}
  \label{fig:ch4-k6-throughput}
\end{figure}

\begin{figure}[!htbp]
  \centering
  \includegraphics[width=0.95\textwidth]{fig03_k6_latency_percentiles.png}
  \caption{Các phân vị độ trễ HTTP theo kịch bản k6}
  \label{fig:ch4-k6-latency}
\end{figure}

Ba kịch bản tải gồm truy vấn đọc dữ liệu, gợi ý phản biện viên và kiểm tra xung đột lợi ích. Mỗi kịch bản sử dụng 20 người dùng ảo trong 30 giây. Kịch bản được gọi là CRUD chỉ thực hiện ba thao tác đọc: liệt kê hội nghị, liệt kê bài nộp và tìm kiếm người dùng; kịch bản không tạo, cập nhật hoặc xóa bản ghi. Cấu hình đầy đủ được trình bày tại Mục~\ref{appc:backend-contract}.

\begingroup
\setlength{\tabcolsep}{4pt}
\begin{longtable}{|>{\RaggedRight\arraybackslash}p{0.19\textwidth}|>{\centering\arraybackslash}p{0.13\textwidth}|>{\centering\arraybackslash}p{0.16\textwidth}|>{\centering\arraybackslash}p{0.14\textwidth}|>{\centering\arraybackslash}p{0.12\textwidth}|>{\centering\arraybackslash}p{0.12\textwidth}|}
  \caption{Kết quả tải HTTP theo kịch bản k6}\label{tab:ch4-demo-http}\\[0.5\baselineskip]
  \hline
  Kịch bản & Số yêu cầu & Thông lượng & Độ trễ trung vị & Độ trễ p95 & Tỷ lệ thất bại \\ \hline
  \endfirsthead
  \hline
  Kịch bản & Số yêu cầu & Thông lượng & Độ trễ trung vị & Độ trễ p95 & Tỷ lệ thất bại \\ \hline
  \endhead
  \hline
  \endfoot
  Truy vấn đọc & 11.110 & 369 yêu cầu/giây & 46,2 ms & 117,6 ms & 0,00\% \\ \hline
  Đối sánh phản biện viên & 17.184 & 572 yêu cầu/giây & 9,7 ms & 71,8 ms & 0,00\% \\ \hline
  Xung đột lợi ích & 16.760 & 558 yêu cầu/giây & 9,5 ms & 79,3 ms & 0,00\% \\ \hline
\end{longtable}
\endgroup

Bảng~\ref{tab:ch4-demo-http} cho thấy backend xử lý hàng trăm yêu cầu mỗi giây và giữ độ trễ p95 dưới 120 ms trong cả ba kịch bản. Kết quả chỉ phản ánh tải ngắn hạn trên các đường xử lý đã chọn; phép thử không thay thế kiểm thử đầy đủ vòng đời chức năng, kiểm thử chạy bền hoặc đánh giá tải phân tán.

\subsection{Tài nguyên tiêu thụ}

Container API sử dụng trung bình 28\% công suất của một nhân CPU, cao nhất 43\%, và khoảng 30 MB RAM. PostgreSQL sử dụng trung bình khoảng 115\% công suất của một nhân CPU, cao nhất 163\%, với mức RAM từ 204 đến 222 MB. Neo4j dùng khoảng 508 MB RAM nhưng dưới 1\% CPU, còn Redis dùng khoảng 9 MB RAM. Hình~\ref{fig:ch4-resource-cpu} và Hình~\ref{fig:ch4-resource-ram} cho thấy cùng xu hướng theo thành phần. Vì CPU được báo cáo theo một nhân và phép đo chưa tách thời gian truy vấn theo thành phần, kết quả chỉ xác định PostgreSQL là thành phần cần ưu tiên kiểm tra khi tối ưu, không chứng minh đây là điểm nghẽn duy nhất ở cấu hình khác.

\begin{figure}[!htbp]
  \centering
  \includegraphics[width=0.94\textwidth]{fig09_resource_cpu.png}
  \caption{Mức sử dụng CPU theo giai đoạn và thành phần}
  \label{fig:ch4-resource-cpu}
\end{figure}

\begin{figure}[!htbp]
  \centering
  \includegraphics[width=0.88\textwidth]{fig10_resource_ram.png}
  \caption{Mức sử dụng RAM theo thành phần}
  \label{fig:ch4-resource-ram}
\end{figure}

\section{Đánh giá lớp thuật toán có thể kiểm chứng}
\label{sec:chap4-deterministic-algo}

\subsection{Chi phí xử lý}

Hệ thống tính độ tương đồng Jaccard theo công thức $J(A,B)=|A\cap B|/|A\cup B|$~\cite{ch4ref7}, sau đó áp dụng thủ tục gán tham lam có xét giới hạn tải. Go microbenchmark đo trực tiếp mã thuật toán và không bao gồm HTTP, cơ sở dữ liệu hoặc mạng. Hình~\ref{fig:ch4-microbench-scale} trình bày thời gian xử lý theo quy mô đầu vào, còn Hình~\ref{fig:ch4-microbench-allocs} trình bày số lần cấp phát bộ nhớ. Hai hình thay thế cho bảng số liệu chi tiết để làm nổi bật xu hướng tăng chi phí khi quy mô đầu vào mở rộng.

\begin{figure}[!htbp]
  \centering
  \includegraphics[width=0.90\textwidth]{fig07_microbench_scale.png}
  \caption{Thời gian Go microbenchmark theo quy mô đầu vào}
  \label{fig:ch4-microbench-scale}
\end{figure}

\begin{figure}[!htbp]
  \centering
  \includegraphics[width=0.90\textwidth]{fig08_microbench_allocs.png}
  \caption{Số lần cấp phát bộ nhớ trong Go microbenchmark}
  \label{fig:ch4-microbench-allocs}
\end{figure}

Chi phí đối sánh tăng từ 131 $\mu$s lên 56 ms khi ma trận điểm giữa bài và hồ sơ mở rộng. Chi phí của hai bộ kiểm tra xung đột xác định tăng từ 14,9 $\mu$s lên 653 $\mu$s. Các giá trị này xác nhận chi phí mã thuật toán trong các quy mô đã thử, nhưng không đại diện cho độ trễ đầu cuối hoặc lớp đồ thị Neo4j.

\subsection{Hành vi xếp hạng và phân công}

Phép thử leave-one-out giữ lại một bài từ mỗi hồ sơ tác giả làm truy vấn, sau đó xếp hạng 60 hồ sơ bằng Jaccard, số chủ đề chung và phương pháp ngẫu nhiên. Nhãn tác giả gốc chỉ kiểm tra khả năng truy hồi quan hệ chủ đề đã được mã hóa trong dữ liệu; nhãn này không đại diện cho phản biện viên phù hợp vì tác giả của bài còn là một xung đột lợi ích hiển nhiên trong vận hành thật.

\begingroup
\setlength{\tabcolsep}{3.5pt}
\begin{longtable}{|>{\RaggedRight\arraybackslash}p{0.21\textwidth}|>{\centering\arraybackslash}p{0.13\textwidth}|>{\centering\arraybackslash}p{0.13\textwidth}|>{\centering\arraybackslash}p{0.13\textwidth}|>{\centering\arraybackslash}p{0.13\textwidth}|>{\centering\arraybackslash}p{0.15\textwidth}|}
  \caption{Kết quả xếp hạng hồ sơ tác giả}\label{tab:ch4-demo-ranking}\\[0.5\baselineskip]
  \hline
  Phương pháp & Hit@1 & Hit@5 & Hit@10 & MRR & nDCG@10 \\ \hline
  \endfirsthead
  \hline
  Phương pháp & Hit@1 & Hit@5 & Hit@10 & MRR & nDCG@10 \\ \hline
  \endhead
  \hline
  \endfoot
  Jaccard & 0,250 & 0,550 & 0,650 & 0,392 & 0,442 \\ \hline
  Số chủ đề chung & 0,233 & 0,550 & 0,733 & 0,391 & 0,463 \\ \hline
  Ngẫu nhiên & 0,017 & 0,083 & 0,167 & 0,078 & 0,076 \\ \hline
\end{longtable}
\endgroup

Hình~\ref{fig:ch4-ranking-hit} cho thấy Jaccard và phương pháp đếm số chủ đề chung đều vượt đường cơ sở ngẫu nhiên ở các mức Hit@k. Jaccard đạt MRR 0,392, cao gấp khoảng năm lần phương pháp ngẫu nhiên. Tuy nhiên, phương pháp đếm số chủ đề chung đạt Hit@10 và nDCG@10 cao hơn Jaccard. Kết quả trên 60 truy vấn chưa cho thấy một phương pháp chiếm ưu thế ổn định.

\begin{figure}[!htbp]
  \centering
  \includegraphics[width=0.82\textwidth]{fig02_ranking_hit_at_k.png}
  \caption{Tỷ lệ Hit@k theo phương pháp xếp hạng}
  \label{fig:ch4-ranking-hit}
\end{figure}

Vì chưa có phương án phân công tối ưu do Chủ tọa xác nhận, nhóm đã so sánh thuật toán Greedy với gán tuần tự và ngẫu nhiên bằng các chỉ số nội tại. Greedy đạt điểm Jaccard trung bình 0,011, cao gấp 2,75 lần hai phương pháp cơ sở, nhưng độ phủ đủ hai phản biện viên chỉ đạt 65,9\%. Có 23,3\% số bài cần lượt dự phòng; lượt này bỏ ngưỡng điểm và giới hạn tải nhưng vẫn giữ xung đột lợi ích là ràng buộc cứng. Hình~\ref{fig:ch4-assignment-tradeoff} minh họa đánh đổi giữa độ phủ và điểm phù hợp trung bình. Kết quả chỉ hỗ trợ sử dụng Greedy để tạo danh sách ứng viên cho Chủ tọa kiểm tra, không hỗ trợ phân công tự động.

\begin{figure}[!htbp]
  \centering
  \includegraphics[width=0.84\textwidth]{fig06_assignment_coverage_score.png}
  \caption{Đánh đổi giữa độ phủ và điểm Jaccard trung bình của phân công}
  \label{fig:ch4-assignment-tradeoff}
\end{figure}

\subsection{Phát hiện xung đột lợi ích đa tầng}

Cơ chế phát hiện xung đột lợi ích gồm kiểm tra tự gán, quan hệ do người dùng khai báo và quan hệ đồng tác giả từ một đến ba bậc được truy vấn theo mẫu đường đi trên Neo4j~\cite{ch4ref6} trong khoảng thời gian cấu hình. Microbenchmark chỉ bao phủ hai bộ kiểm tra xác định đầu tiên. Phép đánh giá chưa đo thời gian hoặc chất lượng phát hiện của lớp đồ thị, chưa có tập cặp xung đột do chuyên gia gán nhãn và chưa kiểm tra đầu cuối khi Neo4j không khả dụng. Vì vậy, kết quả hiện tại chưa xác nhận chất lượng của toàn bộ cơ chế đa tầng.

\section{Đánh giá các luồng AI hỗ trợ}
\label{sec:chap4-ai-flows}

\subsection{Submission Autofill và Submission Gating}

Trên 1.127 bài, chức năng Tự động điền bài nộp (Submission Autofill) đạt F1 theo token của tiêu đề 98,20\%, F1 từ khóa 92,77\% và F1 tác giả 83,49\%. Tỷ lệ hoàn tất trường bắt buộc đạt 86,93\%. Các trường ngắn và có cấu trúc rõ đạt mức khớp cao hơn trường có cấu trúc phức tạp. Hình~\ref{fig:ch4-autofill-quality} minh họa sự khác biệt giữa các nhóm trường. Giá trị thấp nhất bằng 0 ở một số trường cho thấy luồng có những trường hợp thất bại hoàn toàn; hệ thống vì vậy chỉ nên tạo bản nháp để Tác giả kiểm tra và chỉnh sửa.

\begin{figure}[H]
  \centering
  \includegraphics[width=0.94\textwidth]{fig08_autofill_quality.png}
  \caption{Mức khớp của các trường siêu dữ liệu trong Submission Autofill}
  \label{fig:ch4-autofill-quality}
\end{figure}

Track Recommendation trong Submission Autofill nhận nội dung bài cùng danh sách chuyên đề hợp lệ của hội nghị, sau đó trả về các chuyên đề đề xuất trong chính danh sách này. Bộ kiểm thử riêng gồm 48 trường hợp; luồng hoàn tất 48/48 trường hợp, không tạo đề xuất ngoài danh sách hợp lệ và có độ trễ trung bình 18,19 giây. Kết quả xác nhận tính hợp lệ của đầu ra và khả năng hoàn tất luồng. Tập kiểm thử chưa có nhãn chuyên gia về chuyên đề phù hợp nhất, nên kết quả không được dùng để kết luận độ chính xác Top-1 hoặc Top-3.

Submission Gating tách hai tuyến đánh giá. Tuyến kiểm tra luật cố định đối chiếu trực tiếp tệp và siêu dữ liệu với các điều kiện nộp bài có kết quả kỳ vọng; tuyến LLM Gating đọc nội dung để tạo cảnh báo không chặn đối với các vấn đề cần Tác giả hoặc người phụ trách xem xét. Nhánh kiểm tra luật hoàn tất 8/8 trường hợp, đạt 100\% về phán quyết và mã luật, đồng thời không ghi nhận trường hợp chặn sai. Nhánh LLM Gating hoàn tất 24/24 trường hợp, tạo 26 cảnh báo và không có trường hợp biến cảnh báo nội dung thành quyết định chặn tự động. Kết quả này xác nhận luồng giữ đúng ranh giới hỗ trợ, nhưng chưa xác nhận độ chính xác hoặc khả năng hỗ trợ chỉnh sửa của từng cảnh báo vì tập kiểm thử chưa có nhãn chuyên gia tương ứng. Cấu trúc của các bộ kiểm thử được trình bày tại Mục~\ref{appc:direct-fixtures}.

\subsection{Reviewer Initial Analysis}

Thử nghiệm cung cấp bản thảo và siêu dữ liệu bài nộp cho Reviewer Initial Analysis. Luồng tạo các trích dẫn, chú giải và điểm cần lưu ý nhằm định hướng Phản biện viên trước khi viết nhận xét. TCA kiểm tra riêng mức bám nguồn của trích dẫn và các nhận định mở; các bản phản biện của con người được sử dụng để đo Coverage và Additionality, không được xem là đáp án duy nhất của bài toán.

Chức năng Phân tích và tóm tắt bài báo (Reviewer Initial Analysis) tạo bản định hướng trước khi Phản biện viên viết nhận xét. Theo phép chấm TCA, tỷ lệ trích dẫn bám nguồn đạt 96,22\%, tỷ lệ nội dung bị phân loại là không bám nguồn đạt 3,78\%, còn Truthfulness của điểm cần lưu ý đạt 69,86\%. Hình~\ref{fig:ch4-reviewer-grounding} tách mức bám nguồn của trích dẫn khỏi các điểm cần lưu ý. Coverage đạt 4,49\% và Additionality đạt 92,23\%; hai chỉ số này không trực tiếp đo mức độ mới hoặc hữu ích của nội dung bổ sung.

Kết quả cho thấy luồng có thể hỗ trợ định hướng đọc trong phạm vi phép chấm, nhưng Phản biện viên vẫn phải kiểm tra từng nhận định trên bài gốc. Vì đầu ra có thể ảnh hưởng đến phán đoán chuyên môn, mức độ phù hợp với chính sách sử dụng AI phải được xem xét theo từng hội nghị.

\begin{figure}[H]
  \centering
  \includegraphics[width=0.94\textwidth]{fig09a_reviewer_annotation_grounding.png}
  \caption{Mức bám nguồn của đầu ra Reviewer Initial Analysis theo phép chấm TCA}
  \label{fig:ch4-reviewer-grounding}
\end{figure}

\subsection{Review Quality Auditor}

Thử nghiệm cung cấp nội dung bản phản biện cho Review Quality Auditor và ghi nhận danh sách phát hiện, trạng thái cùng mức độ nghiêm trọng. TCA kiểm tra liệu mỗi phát hiện có được nội dung bản phản biện hỗ trợ và có phù hợp với tiêu chí kiểm tra hay không. Phép đánh giá không xem trạng thái \CodeBreak{block} là một quyết định từ chối thao tác, vì trạng thái này chỉ làm nổi bật nội dung cần người dùng xem xét.

Trên 1.127 bài, chức năng Kiểm tra chất lượng nhận xét (Review Quality Auditor) tạo 3.658 lượt kiểm tra, gồm 1.913 lượt mang trạng thái \textit{block}, 1.650 lượt \textit{warn} và 95 lượt \textit{pass}. Trong luồng hiện hành, \textit{severity=blocking} biểu thị mức nghiêm trọng của phát hiện, còn \textit{status=block} làm nổi bật bản phản biện cần được xem xét trên giao diện. Cả hai nhãn chỉ phục vụ hiển thị và hỗ trợ Phản biện viên chỉnh sửa, không ngăn thao tác nộp bản phản biện. Vì vậy, các số liệu này phản ánh mức độ ưu tiên của phát hiện, không phải số lần hệ thống từ chối thao tác.

Theo TCA, Truthfulness so với bản phản biện đạt 58,28\%, tỷ lệ hợp lệ đạt 71,04\%, còn tỷ lệ phát hiện vừa bám nguồn vừa hợp lệ đạt 46,99\%. Hình~\ref{fig:ch4-quality-auditor} cho thấy sự khác biệt giữa tính hợp lệ, mức bám nguồn và phần đồng thời đạt cả hai điều kiện. Kết quả cho thấy Review Quality Auditor phù hợp với vai trò hiển thị các điểm cần kiểm tra trước khi nộp. Phản biện viên vẫn quyết định có chỉnh sửa bản phản biện theo các phát hiện này hay không và chịu trách nhiệm đối với nội dung được nộp.

\begin{figure}[H]
  \centering
  \includegraphics[width=0.94\textwidth]{fig09c_review_quality_auditor.png}
  \caption{Các chỉ số bám nguồn và tính hợp lệ của Review Quality Auditor}
  \label{fig:ch4-quality-auditor}
\end{figure}

\subsection{Chair Decision Copilot}

Thử nghiệm cung cấp cho Chair Decision Copilot các bản phản biện, phản hồi của tác giả và ngữ cảnh nguồn tương ứng. Luồng tạo cơ sở bằng chứng cùng bản tổng hợp các điểm đồng thuận và bất đồng để Chủ tọa xem xét. TCA đo mức bám nguồn của các nhận định và quan hệ với nội dung tham chiếu; phép đánh giá không yêu cầu hoặc chấm một nhãn chấp nhận hay từ chối bài.

Chair Decision Copilot tổng hợp nhận xét phản biện, phản hồi của tác giả và các điểm đồng thuận hoặc bất đồng trước khi Chủ tọa ra quyết định. Theo TCA, Truthfulness của cơ sở bằng chứng đạt 87,34\%, còn Truthfulness của bản tổng hợp bất đồng đạt 87,11\%. Hình~\ref{fig:ch4-chair-evidence} minh họa các chỉ số của cơ sở bằng chứng và bản tổng hợp bất đồng; hình chỉ mô tả kết quả TCA, không phải độ chính xác của quyết định học thuật.

\begin{figure}[H]
  \centering
  \includegraphics[width=0.94\textwidth]{fig09d_chair_evidence_basis.png}
  \caption{Mức bám nguồn của cơ sở bằng chứng và bản tổng hợp bất đồng}
  \label{fig:ch4-chair-evidence}
\end{figure} Coverage của hai nhóm lần lượt đạt 5,27\% và 13,82\%; Additionality của cơ sở bằng chứng đạt 91,63\%. Tỷ lệ rủi ro cao đạt 1,28\%, tương đương khoảng 14 bài trong tập chấm.

Các kết quả chỉ xác nhận mức bám nguồn của bản tổng hợp theo phép chấm TCA. Quá trình kiểm thử chưa đo mức độ đầy đủ và hữu ích theo đánh giá của Chủ tọa, thời gian đọc, độ chính xác quyết định hoặc thiên lệch tự động hóa. Chủ tọa vì vậy vẫn phải kiểm tra nguồn và chịu trách nhiệm về quyết định cuối cùng.

\subsection{Chatbot Agent}
\label{subsec:chap4-chatbot}

Chatbot Agent được đánh giá như một trợ lý truy vấn dữ liệu ConferenceSpace theo quyền của người dùng, không phải trợ lý phân tích nội dung học thuật hoặc hệ thống trò chuyện tự do. Bộ thử gồm tám nhóm kịch bản, mỗi nhóm có năm biến thể diễn đạt. Trợ lý hoàn tất luồng phản hồi trong cả 40 hội thoại; báo cáo rà soát thủ công hồi cứu phân loại 25 lượt đạt, 12 lượt đạt một phần và 3 lượt không đạt. Tỷ lệ gọi công cụ thành công đạt 97/128 lượt, tương ứng 75,78\%. Hình~\ref{fig:ch4-chatbot-outcomes} và Hình~\ref{fig:ch4-chatbot-tools} tách kết quả hội thoại khỏi khả năng gọi công cụ; tỷ lệ gọi công cụ không được diễn giải như tỷ lệ hoàn tất hội thoại.

\begin{figure}[!htbp]
  \centering
  \includegraphics[width=0.78\textwidth]{fig07_chatbot_manual_outcomes.png}
  \caption{Phân bố kết quả rà soát thủ công của Chatbot Agent}
  \label{fig:ch4-chatbot-outcomes}
\end{figure}

\begin{figure}[!htbp]
  \centering
  \includegraphics[width=0.94\textwidth]{fig06_chatbot_tool_success.png}
  \caption{Kết quả gọi công cụ của Chatbot Agent}
  \label{fig:ch4-chatbot-tools}
\end{figure}

Hình~\ref{fig:ch4-chatbot-response-stages} phân tách thời gian phản hồi của Chatbot Agent thành thời gian đến token đầu tiên, thời gian đến token đầu tiên của câu trả lời hoàn chỉnh, thời gian truyền nội dung và tổng thời lượng. Thời gian đến token đầu tiên đạt trung bình 2,36 giây, nhưng câu trả lời hoàn chỉnh chỉ bắt đầu sau trung bình 23,02 giây do luồng còn thực hiện các lượt gọi công cụ và tổng hợp dữ liệu. Tổng thời lượng trung bình đạt 26,53 giây; tỷ lệ gọi công cụ thành công đạt 75,78\%.

Thử nghiệm không ghi nhận rò rỉ dữ liệu riêng tư trong kịch bản quyền. Tuy nhiên, một biến thể ngoài phạm vi vẫn tạo báo cáo thị trường và một số câu trả lời trình bày lỗi truy vấn thay vì giải thích rõ ranh giới quyền. Các trường chấm theo từng lượt chưa được điền đầy đủ, đồng thời tư liệu chưa ghi người chấm và ngày chấm; vì vậy, phân bố 25/12/3 chỉ là bằng chứng mô tả. Ma trận kịch bản và yêu cầu truy vết được trình bày tại Mục~\ref{appc:chatbot}.

\begin{figure}[H]
  \centering
  \includegraphics[width=0.94\textwidth]{fig14_chatbot_response_stages.png}
  \caption{Các giai đoạn phản hồi và kết quả gọi công cụ của Chatbot Agent}
  \label{fig:ch4-chatbot-response-stages}
\end{figure}

\section{Tính khả thi vận hành}

Các luồng xử lý được đo theo những đơn vị khác nhau: Submission Autofill, Reviewer Initial Analysis và Chair Decision Copilot được tính theo bài; Review Quality Auditor được tính theo lượt kiểm tra; Chatbot Agent được tính theo hội thoại. Hình~\ref{fig:ch4-resource-latency} trình bày độ trễ trung bình, trung vị và cao nhất theo từng luồng. Hình~\ref{fig:ch4-resource-tokens} trình bày token trung bình trên mỗi đơn vị công việc. Do mẫu số khác nhau, số token giữa các luồng không được dùng để xếp hạng chi phí.

\begin{figure}[!htbp]
  \centering
  \includegraphics[width=0.94\textwidth]{fig11_resource_latency.png}
  \caption{Độ trễ theo luồng xử lý AI}
  \label{fig:ch4-resource-latency}
\end{figure}

\begin{figure}[!htbp]
  \centering
  \includegraphics[width=0.94\textwidth]{fig12_resource_tokens.png}
  \caption{Token trung bình theo đơn vị công việc của các luồng AI}
  \label{fig:ch4-resource-tokens}
\end{figure}

Hình~\ref{fig:ch4-resource-ops-mode} chuyển các mức độ trễ quan sát được thành khuyến nghị về cách tổ chức vận hành. Đây là khuyến nghị thiết kế dựa trên độ trễ và đặc điểm phản hồi của từng luồng, không phải kết quả so sánh giữa các phương án triển khai thực tế. Nhánh kiểm tra luật của chức năng Kiểm tra trước khi nộp bài (Submission Gating) có thể chạy đồng bộ. Nhánh LLM Gating phù hợp với xử lý song song không chặn thao tác. Reviewer Initial Analysis, Review Quality Auditor và Chair Decision Copilot có các trường hợp vượt 100 giây, nên phù hợp hơn với cơ chế chạy nền hoặc chạy trước. Chatbot Agent cần phản hồi từng phần và hiển thị trạng thái tra cứu vì thời gian đến token đầu tiên chỉ đạt 2,36 giây, trong khi câu trả lời hoàn chỉnh bắt đầu sau trung bình 23,02 giây.

\begin{figure}[!htbp]
  \centering
  \includegraphics[width=0.94\textwidth]{fig13_resource_ops_mode.png}
  \caption{Khuyến nghị cách tổ chức vận hành theo luồng AI}
  \label{fig:ch4-resource-ops-mode}
\end{figure}

Các kết quả dẫn đến ba yêu cầu vận hành: chuyển tác vụ dài sang hàng đợi nền; hiển thị tiến độ, lỗi và khả năng thử lại; giảm 31 lượt gọi công cụ thất bại của Chatbot Agent.

\section{Khảo sát người dùng}
\label{sec:chap4-uat}

Báo cáo tổng hợp ghi nhận 91 phản hồi dùng cho phân tích, gồm 76 phản hồi của Tác giả, 7 phản hồi của Phản biện viên và 8 phản hồi của Chủ tọa. Đây là mẫu thuận tiện và lệch mạnh về nhóm Tác giả; kết quả chỉ mô tả cảm nhận của mẫu tham gia. Riêng số liệu Chủ tọa chưa thể được tái lập độc lập từ tệp dữ liệu thô hiện có, nên được xem là số liệu tổng hợp thứ cấp. Bản ghi thô đang lưu chỉ có sáu bản ghi Chủ tọa, trong đó năm bản ghi thiếu phần đánh giá chính; vì vậy, các thống kê với mẫu số tám được giữ theo báo cáo tổng hợp nhưng không được xem là đã xác minh độc lập.

Bảng~\ref{tab:ch4-demo-uat} tổng hợp kết quả khảo sát theo ba vai trò.

\begingroup
\setlength{\tabcolsep}{4pt}
\begin{longtable}{|>{\RaggedRight\arraybackslash}p{0.18\textwidth}|>{\centering\arraybackslash}p{0.13\textwidth}|>{\centering\arraybackslash}p{0.22\textwidth}|>{\RaggedRight\arraybackslash}p{0.39\textwidth}|}
  \caption{Tổng hợp kết quả khảo sát theo vai trò}\label{tab:ch4-demo-uat}\\[0.5\baselineskip]
  \hline
  Vai trò & Cỡ mẫu & Hài lòng tổng thể & Tín hiệu chính \\ \hline
  \endfirsthead
  \hline
  Vai trò & Cỡ mẫu & Hài lòng tổng thể & Tín hiệu chính \\ \hline
  \endhead
  \hline
  \endfoot
  Tác giả & 76 & 3,89/5; 67/76 tích cực & 47/76 chọn Submission Autofill là tính năng hữu ích nhất; 69/76 cảm nhận AI giúp giảm thao tác hoặc thời gian \\ \hline
  Phản biện viên & 7 & 4,29/5 & 6/7 phản hồi tích cực với nhóm chức năng hỗ trợ đọc và phản biện \\ \hline
  Chủ tọa & 8 & 4,38/5; 7/8 tích cực & 8/8 phản hồi tích cực về theo dõi tiến độ; 7/8 phản hồi tích cực trong câu hỏi gộp về nhóm công cụ AI \\ \hline
\end{longtable}
\endgroup

Ở nhóm Tác giả, chức năng Tự động điền bài nộp (Submission Autofill) nhận phản hồi tích cực nhất, phù hợp với kết quả đối sánh siêu dữ liệu. Tuy nhiên, khảo sát chỉ đo cảm nhận giảm thao tác hoặc tiết kiệm thời gian, không đo số phút thực tế. Ở nhóm Phản biện viên, câu hỏi gộp nhiều thành phần nên không thể quy kết quả riêng cho chức năng Phân tích và tóm tắt bài báo (Reviewer Initial Analysis). Ở nhóm Chủ tọa, mức độ sử dụng các tính năng không đồng đều và câu hỏi dùng mẫu số theo từng mục; kết quả vì vậy không đánh giá riêng chức năng Báo cáo tổng hợp hỗ trợ Chủ tọa (Chair Decision Copilot) hoặc chất lượng quyết định.

Toàn bộ mẫu có 73/91 người tham gia cho biết sẵn sàng giới thiệu nền tảng. Tỷ lệ này chịu ảnh hưởng lớn từ nhóm Tác giả, vốn chiếm 83,5\% mẫu. UAT cung cấp bằng chứng cảm nhận ban đầu về tính hữu ích, đồng thời nhấn mạnh nhu cầu giữ quyền kiểm tra, sửa và bỏ qua đầu ra AI. Phụ lục~\ref{appc:uat} trình bày nguồn gốc và giới hạn truy vết của dữ liệu khảo sát.

\section{Tổng hợp kết quả và giới hạn}

Bảng~\ref{tab:ch4-demo-synthesis} tổng hợp mức kết luận cho từng nhóm trách nhiệm và bằng chứng tương ứng.

\begingroup
\setlength{\tabcolsep}{4pt}
\renewcommand{\arraystretch}{1.3}
\begin{longtable}{|>{\raggedright\arraybackslash}p{0.27\textwidth}|>{\raggedright\arraybackslash}p{0.43\textwidth}|>{\raggedright\arraybackslash}p{0.22\textwidth}|}
\caption{Mức kết luận theo trách nhiệm và bằng chứng}
\label{tab:ch4-demo-synthesis} \\
\hline
\TableHeader{Nội dung đánh giá} & \TableHeader{Bằng chứng chính} & \TableHeader{Mức kết luận} \\ \hline
\endfirsthead
\hline
\TableHeader{Nội dung đánh giá} & \TableHeader{Bằng chứng chính} & \TableHeader{Mức kết luận} \\ \hline
\endhead
\hline
\multicolumn{3}{|r|}{\textit{(Tiếp tục ở trang sau...)}} \\ \hline
\endfoot
\hline
\endlastfoot
Hiệu năng của các endpoint đã chọn & Ba kịch bản k6; độ trễ p95 dưới 120 ms; tỷ lệ yêu cầu thất bại bằng 0\% & Bằng chứng trực tiếp trong cấu hình thử nghiệm \\ \hline
Đối sánh và phân công & Phép thử leave-one-out và các phương pháp cơ sở trên dữ liệu tổng hợp; MRR của Jaccard đạt 0,392; độ phủ của Greedy đạt 65,9\% & Chỉ số gián tiếp \\ \hline
Xung đột lợi ích & Microbenchmark cho kiểm tra tự gán và quan hệ khai báo & Chưa đo chất lượng đa tầng \\ \hline
Submission Autofill và luật của Submission Gating & Đối chiếu trên 1.127 bài; 8/8 trường hợp kiểm thử luật & Bằng chứng trực tiếp trong phạm vi đầu vào và tập kiểm thử \\ \hline
Reviewer Initial Analysis & TCA trên 1.097 bài có kết quả đủ điều kiện; tỷ lệ trích dẫn bám nguồn 96,22\%; Truthfulness của điểm cần lưu ý 69,86\% & Chỉ số gián tiếp \\ \hline
Review Quality Auditor & 3.658 lượt; tỷ lệ phát hiện vừa bám nguồn vừa hợp lệ 46,99\% & Phát hiện được hiển thị để hỗ trợ Phản biện viên kiểm tra \\ \hline
Chair Decision Copilot & Truthfulness của cơ sở bằng chứng và bản tổng hợp bất đồng lần lượt đạt 87,34\% và 87,11\% & Chỉ số gián tiếp \\ \hline
Chatbot Agent & 40 hội thoại; 97/128 lượt gọi công cụ thành công; không ghi nhận rò rỉ dữ liệu trong kịch bản quyền & Bằng chứng theo kịch bản \\ \hline
Trải nghiệm người dùng & 91 phản hồi gồm 76 Tác giả, 7 Phản biện viên và 8 Chủ tọa & Bằng chứng cảm nhận \\
\end{longtable}
\endgroup

Kết quả kỹ thuật và UAT cùng cho thấy chức năng Tự động điền bài nộp (Submission Autofill) có thể hỗ trợ nhập liệu, nhưng quá trình kiểm thử chưa đo trực tiếp số thao tác hoặc thời gian tiết kiệm được. Các phép truy hồi chủ đề và TCA cung cấp chỉ số gián tiếp; chúng không xác nhận mức độ phù hợp của phản biện viên, chất lượng quyết định học thuật hoặc tính hữu ích theo đánh giá của chuyên gia. Với chức năng Kiểm tra chất lượng nhận xét (Review Quality Auditor), các phát hiện được hiển thị để Phản biện viên xem xét và chỉnh sửa bản phản biện; trạng thái \textit{block} không ngăn thao tác nộp. Tỷ lệ phát hiện vừa bám nguồn vừa hợp lệ đạt 46,99\%, nên Phản biện viên vẫn phải kiểm tra nội dung phát hiện trước khi sử dụng.

Các giới hạn chính gồm thiếu nhãn chuyên gia cho một số luồng, dữ liệu đối sánh tổng hợp, phạm vi bài báo chủ yếu bằng tiếng Anh, thiếu thử nghiệm chạy bền và thử nghiệm chủ động gây lỗi, cùng khả năng truy vết chưa đầy đủ của UAT ở nhóm Chủ tọa. Bộ thực thi đạt tỷ lệ hoàn tất 97,34\%; 30 bài không hoàn tất do lỗi dịch vụ AI hoặc lỗi trong quá trình chạy benchmark được giữ trong thống kê vận hành nhưng không được đưa vào phép chấm TCA. Cơ chế phân quyền cấp ứng dụng cũng chưa cung cấp bằng chứng về chính sách lưu giữ, sử dụng và khu vực lưu trữ dữ liệu của nhà cung cấp AI.

Tóm lại, Chương 4 xác nhận hiệu năng của các endpoint đã chọn, mức độ khớp của siêu dữ liệu do chức năng Tự động điền bài nộp (Submission Autofill) tạo và hành vi của bộ kiểm thử luật trong cấu hình thử nghiệm. Đánh giá đối sánh, TCA, Chatbot Agent và UAT cung cấp bằng chứng theo phạm vi hẹp hơn. Các kết quả chưa xác nhận chất lượng của toàn bộ cơ chế phát hiện xung đột lợi ích đa tầng, độ chính xác của quyết định học thuật hoặc mức độ sẵn sàng triển khai trong một hội nghị thực tế.
